% Grado 2 -- generado por tools/build_tex.py a partir de raw/grado_02.txt
\section{Grado 2}\label{grado:2}

% PDF p. 15
\dba{1}{Interpreta, propone y resuelve problemas aditivos (de composición, transformación y relación) que involucren la cantidad en una colección, la medida de magnitudes (longitud, peso, capacidad y duración de eventos) y problemas multiplicativos sencillos.}

\evidencias

\begin{itemize}
  \item Interpreta y construye diagramas para representar relaciones aditivas y multiplicativas entre cantidades que se presentan en situaciones o fenómenos.
  \item Describe y resuelve situaciones variadas con las operaciones de suma y resta en problemas cuya estructura puede ser a + b = ?, a + ? = c, o ? + b = c.
  \item Reconoce en diferentes situaciones relaciones aditivas y multiplicativas y formula problemas a partir de ellas.
\end{itemize}

\ejemplo

En una sala de videojuegos se requiere tener dinero para pagar el valor de cada hora. Con base en esta situación:

\begin{itemize}
  \item Propone una pregunta que se pueda responder con una multiplicación.
  \item Propone una pregunta que se pueda responder con una división.
  \item Si hay promoción en la sala y se hace un descuento por cada dos horas de uso del servicio, ¿se podría aplicar reiteradamente la multiplicación para conocer el valor a pagar?
\end{itemize}

% PDF p. 15
\dba{2}{Utiliza diferentes estrategias para calcular (agrupar, representar elementos en colecciones, etc.) o estimar el resultado de una suma y resta, multiplicación o reparto equitativo.}

\evidencias

\begin{itemize}
  \item Construye representaciones pictóricas y establece relaciones entre las cantidades involucradas en diferentes fenómenos o situaciones.
  \item Usa algoritmos no convencionales para calcular o estimar el resultado de sumas, restas, multiplicaciones y divisiones entre números naturales, los describe y los justifica.
\end{itemize}

\ejemplo

Cuatro estudiantes deciden jugar parqués con una sola ficha y con las siguientes reglas: a) En el primer turno se lanzará un solo dado y cada punto de este permitirá mover la ficha tres casillas. b) En el segundo turno se utilizarán dos dados y cada punto permitirá mover la ficha una casilla. Inicia el lanzamiento quien lidera el juego, luego quien vaya de segundo, de tercero, hasta lanzar quien esté en el último puesto. c) En los siguientes turnos se repite la primera regla y luego la segunda regla hasta terminar el juego. Se lanza en el orden que se definió en la segunda regla.

\begin{itemize}
  \item Si al registrar (en su orden) cada lanzamiento en una tabla los resultados son:
\end{itemize}
\fig{g02_p16_1}{Tabla de puntos por jugador (Andrés, Camilo, Juan José, Carlos) en el lanzamiento 1 y el lanzamiento 2.}

Determina el jugador que lidera el juego hasta el momento. Averigua la cantidad de puntos que debe obtener Camilo en el siguiente lanzamiento para liderar el juego después del tercer turno y propone una regla adicional para que sea Juan José quien lidere el juego después de los dos primeros turnos.

% PDF p. 16
\dba{3}{Utiliza el Sistema de Numeración Decimal para comparar, ordenar y establecer diferentes relaciones entre dos o más secuencias de números con ayuda de diferentes recursos.}

\evidencias

\begin{itemize}
  \item Compara y ordena números de menor a mayor y viceversa a través de recursos como la calculadora, aplicación, material gráfico que represente billetes, diagramas de colecciones, etc.
  \item Propone ejemplos y comunica de forma oral y escrita las condiciones que puede establecer para conservar una relación (mayor que, menor que) cuando se aplican algunas operaciones a ellos.
  \item Reconoce y establece relaciones entre expresiones numéricas (hay más, hay menos, hay la misma cantidad ) y describe el tipo de operaciones que debe realizarse para que a pesar de cambiar los valores numéricos, la relación se conserve.
\end{itemize}

\ejemplo

En la imagen se presenta la tarea que Sara hizo en el tablero. Sin embargo, por accidente un compañero borró parte de lo que Sara había hecho.

\fig{g02_p16_2}{Niña frente a un tablero con cadenas de operaciones: 6 + 3 − 6 = □, 10 + 2 = □, 9 − 3 + 5 = □.}

\begin{itemize}
  \item Escribe algunos números en los espacios vacíos para lograr que se cumpla la relación.
  \item Any dijo que el número que se había borrado en la primera línea era 2. Argumenta si Any tiene o no la razón.
  \item Cuando Margarita vio el tablero quiso ayudar a Sara con la segunda línea del tablero. Ella dijo que sólo podría poner dos números en el espacio borrado. Determina la validez del argumento de Margarita y explica por qué ella dice que hay dos posibilidades. Valora la validez de la afirmación: José dijo que Mario se había equivocado en la tercera línea porque para que se cumpla la igualdad, después del 5 no podía quedar ningún número. Justifica su valoración.
\end{itemize}

% PDF p. 17
\dba{4}{Compara y explica características que se pueden medir, en el proceso de resolución de problemas relativos a longitud, superficie, velocidad, peso o duración de los eventos, entre otros.}

\evidencias

\begin{itemize}
  \item Utiliza instrumentos y unidades de medición apropiados para medir magnitudes diferentes.
  \item Describe los procedimientos necesarios para medir longitudes, superficies, capacidades, pesos de los objetos y la duración de los eventos.
  \item Mide magnitudes con unidades arbitrarias y estandarizadas.
  \item Estima la medida de diferentes magnitudes en situaciones prácticas.
\end{itemize}

\ejemplo

Analiza diferentes situaciones en las que se comparan objetos según magnitudes y describe estrategias para: calcular la distancia recorrida por un auto que se mueve a cierta velocidad constante durante un intervalo de tiempo; calcula o estima la cantidad de tela que se gastaría en un vestido, la longitud de una cinta para cubrir el borde de una mesa; busca longitudes cercanas a un metro o pesos cercanos a un kilogramo e identifica otros objetos que podrían tener esa

\fig{g02_p17_1}{Escena de un parque con estructuras y una figura sobre cuadrícula para describir posiciones y trayectos.}

% PDF p. 17
\dba{5}{Utiliza patrones, unidades e instrumentos convencionales y no convencionales en procesos de medición, cálculo y estimación de magnitudes como longitud, peso, capacidad y tiempo.}

\evidencias

\begin{itemize}
  \item Describe objetos y eventos de acuerdo con atributos medibles: superficie, tiempo, longitud, peso, ángulos.
  \item Realiza mediciones con instrumentos y unidades no convencionales, como pasos, cuadrados o rectángulos, cuartas, metros, entre otros.
  \item Compara eventos según su duración, para ello utiliza relojes convencionales.
\end{itemize}

\ejemplo

Pipe y Lupe salen al mismo tiempo de sus lugares respectivos (cuadrado azul y cuadrado verde), pasan por la zanahoria que tienen más cerca y llegan hasta donde está el conejo. En este recorrido Pipe tarda 30 minutos y Lupe tarda 35 minutos.

Señala la pareja de relojes correspondiente a la hora de llegada de los niños hasta el conejo y explica la respuesta.

\fig{g02_p18_1}{Pares de relojes análogos (situaciones 3 y 4) para leer y comparar horas.}

\fig{g02_p18_2}{Otros pares de relojes análogos para leer y comparar horas.}

% PDF p. 18
\dba{6}{Clasifica, describe y representa objetos del entorno a partir de sus propiedades geométricas para establecer relaciones entre las formas bidimensionales y tridimensionales.}

\evidencias

\begin{itemize}
  \item Reconoce las figuras geométricas según el número de lados.
  \item Diferencia los cuerpos geométricos.
  \item Compara figuras y cuerpos geométricos y establece relaciones y diferencias entre ambos.
\end{itemize}

\ejemplo

La habitación de Andrés se muestra en la siguiente imagen:

\fig{g02_p18_3}{Habitación de un niño con objetos (balón, cama, repisa, ventana) para describir ubicaciones relativas.}

Sofía es la tía de Andrés y se encuentra en otro país, ella quiere hacerse a una idea de la habitación de su sobrino. Andrés escribe una carta a Sofía en la que describe detalladamente la habitación de Andrés y los objetos que hay en ella. Estudia si la descripción que hace Andrés es correcta, justifica su respuesta y propone afirmaciones que la completan.

Querida tía Sofía. Estoy muy feliz con mi nueva habitación y quiero contarte cómo es. Imagínate que tengo una cama de forma rectangular y un armario. También tengo una ventana cuadrada y una pintura en la cabecera de la cama de forma triangular. ¡Tengo muchos juguetes! Dos balones que tienen forma esférica, de fútbol y de baloncesto. También tengo carros y cubos para apilar. No te imaginas lo feliz que me siento con todo lo que tengo en mi habitación. Espero que me envíes una foto para ponerla en el portarretrato. Un abrazo. Tu sobrino, Andrés.

% PDF p. 19
\dba{7}{Describe desplazamientos y referencia la posición de un objeto mediante nociones de horizontalidad, verticalidad, paralelismo y perpendicularidad en la solución de problemas.}

\evidencias

\begin{itemize}
  \item Describe desplazamientos a partir de las posiciones de las líneas.
  \item Representa líneas y reconoce las diferentes posiciones y la relación entre ellas.
  \item En dibujos, objetos o espacios reales, identifica posiciones de objetos, de aristas o líneas que son paralelas, verticales o perpendiculares.
  \item Argumenta las diferencias entre las posiciones de las líneas.
\end{itemize}

\ejemplo

\begin{itemize}
  \item Identifica desplazamientos en lugares determinados que estén en correspondencia con unas normas establecidas. Da indicaciones para llegar a determinado sitio. Para ello utiliza palabras como: vertical, horizontal, paralelo, perpendicular.
  \item En la figura se muestra el mapa de un lugar; indica de manera verbal, escrita o gráfica cómo llegar de la casa a la iglesia sin perderse o desplazándose con cierta condición.
\end{itemize}
\fig{g02_p19_1}{Mapa de un barrio con calles y marcadores de lugares para describir recorridos.}

% PDF p. 19
\dba{8}{Propone e identifica patrones y utiliza propiedades de los números y de las operaciones para calcular valores desconocidos en expresiones aritméticas.}

\evidencias

\begin{itemize}
  \item Establece relaciones de reversibilidad entre la suma y la resta.
  \item Utiliza diferentes procedimientos para calcular un valor desconocido.
\end{itemize}

\ejemplo

Ubica un número de entrada y efectúa las operaciones indicadas en la cadena numérica.

\fig{g02_p19_2}{Máquina de transformación: número de entrada → +5 → □ → +5 → □ → −2 → número de salida.}

\begin{itemize}
  \item Encuentra los valores respectivos de salida cuando los números de entrada son 1, 4 y 7.
  \item Encuentra los valores de entrada para que los números de salida sean 18 y 36.
\end{itemize}

\fig{g02_p19_3}{Segunda máquina de transformación: número de entrada → □ → número de salida, con operaciones por descubrir.}

Verbaliza las propiedades que se usan cuando recorre la cadena de izquierda a derecha y de derecha a izquierda. Construye cadenas equivalentes a la cadena dada inicialmente.

% PDF p. 20
\dba{9}{Opera sobre secuencias numéricas para encontrar números u operaciones faltantes y utiliza las propiedades de las operaciones en contextos escolares o extraescolares.}

\evidencias

\begin{itemize}
  \item Utiliza las propiedades de las operaciones para encontrar números desconocidos en igualdades numéricas.
  \item Utiliza las propiedades de las operaciones para encontrar operaciones faltantes en un proceso de cálculo numérico.
  \item Reconoce que un número puede escribirse de varias maneras equivalentes.
  \item Utiliza ensayo y error para encontrar valores u operaciones desconocidas.
\end{itemize}

\ejemplo

Encuentra todas las parejas de números cuya suma es 12 y todos los resultados que se obtienen al multiplicar los números de cada pareja. Propone una tabla para presentar y relacionar los resultados. Realiza lo mismo con otros números (10, 11, etc.) y compara las parejas obtenidas. Establece procedimientos para encontrar las parejas y construye reglas para saber cuántas parejas se pueden formar según el resultado de la suma. Propone ideas sobre cuál es la pareja en la que el resultado es mayor.

\fig{g02_p20_1}{Niña pensando: a + b = 12, a × b = ?}

% PDF p. 20
\dba{10}{Clasifica y organiza datos, los representa utilizando tablas de conteo, pictogramas con escalas y gráficos de puntos, comunica los resultados obtenidos para responder preguntas sencillas.}

\evidencias

\begin{itemize}
  \item Identifica la equivalencia de fichas u objetos con el valor de la variable.
  \item Organiza los datos en tablas de conteo y en pictogramas con escala (uno a muchos).
  \item Lee la información presentada en tablas de conteo, pictogramas con escala y gráficos de puntos.
  \item Comunica los resultados respondiendo preguntas tales como: ¿cuántos hay en total?, ¿cuántos hay de cada dato?, ¿cuál es el dato que más se repite?, ¿cuál es el dato que menos se repite?
\end{itemize}

\ejemplo

\begin{itemize}
  \item En el colegio se realizan las elecciones a personero. La información de los resultados de las votaciones se presenta en el siguiente gráfico:
\end{itemize}

\fig{g02_p20_2}{Pictograma de votos para personero (cada estrella = 5 votos): candidato 1: 4 estrellas, candidato 2: 7, candidato 3: 9, candidato 4: 3.}

\begin{itemize}
  \item En el informe que se entrega, se afirma que: a) El ganador fue el candidato 2; b) El total de votos fue de 210; c) El candidato ganador obtuvo el doble de votos que el candidato que obtuvo menos votos; d) El candidato 4 obtuvo la mitad de votos que el candidato 2. Escribe un informe en el que se compara la información textual con la presentada en la gráfica para dar argumentos sobre la veracidad de la información presentada.
\end{itemize}

% PDF p. 21
\dba{11}{Explica, a partir de la experiencia, la posibilidad de ocurrencia o no de un evento cotidiano y el resultado lo utiliza para predecir la ocurrencia de otros eventos.}

\evidencias

\begin{itemize}
  \item Diferencia situaciones cotidianas cuyo resultado puede ser incierto de aquellas cuyo resultado es conocido o seguro.
  \item Identifica resultados posibles o imposibles, según corresponda, en una situación cotidiana
  \item Predice la ocurrencia o no de eventos cotidianos basado en sus observaciones.
\end{itemize}

\ejemplo

En el grado 2º deciden jugar al lanzamiento de aviones de papel. Acuerdan que todos lanzan los aviones desde una raya que dibujan cerca del tablero del salón y gana el estudiante que lance el avión más lejos. David afirma que él será siempre el ganador porque ya sabe lanzar aviones de papel. Determina si la afirmación del niño es verdadera o falsa y justifica su respuesta.

